% !PRINTED none
% Abschnitt-opening recto; no folio is printed.
\begin{center}
\gesperrt{Seventeenth Chapter.}

\textbf{\Large Elliptic Functions and Quadratic Forms.}

\rule{0.12\linewidth}{0.4pt}
\end{center}

\textbf{§ 114. Singular periods of the doubly periodic functions.}

If we denote by $\varphi(u)$ a doubly periodic function with the periods $\omega_1$, $\omega_2$, then, if $n$ denotes an integer, $\varphi(nu)$ has the same periods, and upon this rests the multiplication of the elliptic functions which was considered in the sixth chapter. The question now arises whether a multiplier $\mu$ cannot be determined in yet another way, so that $\varphi(\mu u)$ has the periods $\omega_1$, $\omega_2$; that is, whether a multiplication exists also with a non-integer multiplier. This will be the case then and only then, when for a system of integers $a$, $b$, $c$, $\partial$ the equations hold:
\[
\begin{aligned}
\mu\,\omega_1 &= a\,\omega_1 + b\,\omega_2,\\
\mu\,\omega_2 &= c\,\omega_1 + \partial\,\omega_2.
\end{aligned}
\]
If, in order to explore more closely the requirement contained in them, we put:
\[
\omega = \dfrac{\omega_2}{\omega_1}
\]
and assume that $\omega$ has a positive imaginary part, then it follows:
\[\mu=a+b\omega,\]
\weq{1}{\omega=\dfrac{c+\partial\omega}{a+b\omega}.}
So long as $\omega$ is regarded as variable, this equation is only possible if $b=c=0$, $a=\partial$, and $\mu$ is therefore an integer.
